\section{Interpretación de signos y significancia}

\subsection{Consideraciones de interpretación}

La interpretación de los resultados se realiza usando \texttt{BIP} como alternativa de referencia
para las variables comunes al viaje o a la zona. Por lo tanto, para variables como año, franja
horaria, demanda, oferta, Censo, OSM y EOD, un coeficiente positivo de \texttt{QR\_RED} o
\texttt{QR\_OTHER} indica una mayor utilidad relativa de esa alternativa respecto de \texttt{BIP},
manteniendo constantes las demás variables del modelo. Un coeficiente negativo indica una menor
utilidad relativa respecto de \texttt{BIP}.

En el caso de variables alternativa-específicas, como tiempos de viaje, esperas y transbordos, los
coeficientes se interpretan directamente para cada alternativa. Estas variables pueden tomar valores
distintos para \texttt{BIP}, \texttt{QR\_RED} y \texttt{QR\_OTHER} dentro de una misma observación,
por lo que sus coeficientes son identificables para las tres alternativas.

Finalmente, los resultados deben interpretarse como asociaciones condicionales del modelo, no como
efectos causales. En particular, las variables territoriales pueden capturar diferencias de
composición socioespacial, intensidad de uso del sistema, disponibilidad de infraestructura y otros
factores no observados que coexisten en una misma zona. Esta cautela es especialmente relevante para
las variables EOD 2012, que se incorporan como proxies históricos de composición territorial y no
como mediciones contemporáneas directas del período 2024--2025.

\subsection{Variables de viaje}

Las variables de espera presentan los patrones más estables en ambas especificaciones candidatas. El
tiempo de espera inicial (\texttt{T\_ESPERA\_INI}) tiene coeficientes negativos, de magnitud similar
y altamente significativos para las tres alternativas. Esto es consistente con la teoría de elección
discreta: mayores tiempos de espera reducen la utilidad de la alternativa considerada.

El tiempo de espera asociado a transbordos (\texttt{T\_ESPERA\_TRASB}) también presenta coeficientes
negativos para \texttt{BIP} y \texttt{QR\_OTHER}, ambos estadísticamente significativos. En
\texttt{QR\_RED}, en cambio, el coeficiente es cercano a cero y no significativo en ambas
especificaciones. Una lectura sustantiva posible es que los viajes asociados a \texttt{QR\_RED}
presentan menor penalización por espera en transbordo, posiblemente porque el uso de la app entrega
información que permite anticipar la espera o decidir cómo utilizar ese tiempo. Esta hipótesis debe
tratarse como interpretación exploratoria y no como evidencia causal directa.

El número de transbordos (\texttt{N\_TRASB}) tiene signo negativo para las tres alternativas, aunque
su significancia es mayor para \texttt{QR\_OTHER} y \texttt{QR\_RED}. Esto indica que, manteniendo
constantes las demás variables, viajes con más transbordos tienden a reducir la utilidad asociada a
las alternativas de pago, especialmente en las alternativas QR.

Por su parte, el tiempo en vehículo (\texttt{T\_VEH}) presenta coeficientes de magnitud muy pequeña.
El único coeficiente estadísticamente significativo aparece en \texttt{QR\_OTHER}, con signo
positivo. Esta señal debe interpretarse con cautela: más que indicar una preferencia por viajes más
largos, probablemente refleja diferencias de composición entre tipos de viaje, zonas de origen,
distancias recorridas y alternativas disponibles.

\subsection{Variables temporales, demanda y oferta}

La variable \texttt{ANIO\_2025} muestra una asociación positiva con \texttt{QR\_OTHER} y
\texttt{QR\_RED} respecto de \texttt{BIP}. Esto sugiere que, en la muestra analizada, el año 2025 se
asocia con una mayor utilidad relativa de las alternativas QR. Esta señal es consistente con una
posible expansión o consolidación del uso de medios de pago QR entre 2024 y 2025.

Las variables de franja muestran diferencias relevantes entre alternativas. \texttt{LAB\_PM} se
asocia fuertemente con \texttt{QR\_RED} respecto de \texttt{BIP}, mientras que también presenta una
asociación positiva, aunque menor, con \texttt{QR\_OTHER}. \texttt{LAB\_PT} muestra una señal
similar, aunque de menor magnitud. En contraste, \texttt{NO\_LAB} se asocia positivamente con
\texttt{QR\_OTHER}, mientras que la señal para \texttt{QR\_RED} no es estadísticamente robusta.

La variable \texttt{LOG\_DEMAND} presenta coeficientes negativos y significativos para
\texttt{QR\_OTHER} y \texttt{QR\_RED} respecto de \texttt{BIP}. Esto sugiere que zonas y franjas con
mayor intensidad local de viajes tienden a estar relativamente menos asociadas a alternativas QR. Sin
embargo, esta variable debe interpretarse con cautela, porque puede reflejar tanto intensidad de uso
del sistema como composición territorial de la demanda.

En cuanto a oferta, la densidad de paraderos se asocia negativamente con \texttt{QR\_RED} y no
presenta una señal robusta para \texttt{QR\_OTHER}. El número de líneas de bus y de Metro muestra
coeficientes positivos para \texttt{QR\_RED}, y el número de líneas de Metro también es positivo para
\texttt{QR\_OTHER}. Estos resultados sugieren que la relación entre infraestructura operacional y
medio de pago no es homogénea entre alternativas QR, y probablemente captura diferencias entre tipos
de viaje, zonas de origen y estructura de la red.

\subsection{Variables censales}

La proporción de población de 18 años o más con educación universitaria o superior muestra uno de los
patrones más claros de ambos modelos. El coeficiente es positivo y altamente significativo para
\texttt{QR\_RED} respecto de \texttt{BIP}, y también positivo para \texttt{QR\_OTHER}, aunque de
menor magnitud. En términos relativos, esto indica que las zonas con mayor presencia de población
universitaria se asocian con una mayor utilidad relativa de alternativas QR, especialmente
\texttt{QR\_RED}. Este resultado es coherente con la hipótesis de que mayor nivel educacional puede
relacionarse con mayor adopción de herramientas digitales de pago o información.

El hacinamiento presenta coeficientes negativos para las alternativas QR, especialmente para
\texttt{QR\_OTHER}. En \texttt{QR\_RED} la señal también es negativa, aunque menos robusta. Este
patrón sugiere que zonas con mayor hacinamiento se asocian relativamente menos con alternativas QR,
lo que puede reflejar condiciones socioeconómicas y diferencias de acceso o adopción digital.

La proporción de personas con discapacidad muestra una asociación negativa y significativa con
\texttt{QR\_RED}, mientras que la señal para \texttt{QR\_OTHER} no es robusta en las especificaciones
actuales. En este caso, la interpretación debe ser cuidadosa, porque la variable describe el contexto
territorial de origen y no necesariamente una característica individual del usuario que realiza el
viaje.

La proporción de población inmigrante se asocia positivamente con \texttt{QR\_OTHER} y también con
\texttt{QR\_RED}. La magnitud es mayor para \texttt{QR\_OTHER}, aunque ambas señales se mantienen
positivas en los dos modelos. Esto sugiere que zonas con mayor proporción de población inmigrante
están relativamente más asociadas a alternativas QR respecto de \texttt{BIP}.

Las variables censales de composición sociodemográfica agregan dos lecturas complementarias. La proporción de mujeres en la
zona de origen no presenta una señal robusta para \texttt{QR\_OTHER}, pero sí una asociación positiva
y significativa con \texttt{QR\_RED}. Por su parte, la asistencia a educación parvularia muestra una
señal positiva y significativa para \texttt{QR\_RED}, mientras que su asociación con
\texttt{QR\_OTHER} es débil o no robusta. Dado que estas variables son territoriales, no deben
interpretarse como características individuales del usuario, sino como indicadores de composición
sociodemográfica y ciclo de vida de las zonas de origen.

\subsection{Variables OSM e infraestructura}

La densidad de universidades OSM presenta coeficientes negativos y significativos para
\texttt{QR\_OTHER} y \texttt{QR\_RED} respecto de \texttt{BIP}. Este resultado no contradice
necesariamente el efecto positivo de la educación universitaria residencial sobre \texttt{QR\_RED}.
Ambas variables capturan dimensiones distintas: la variable censal de educación universitaria
describe la composición educacional de la población residente, mientras que
\path{osm_amenity_university_density_km2_z} identifica la presencia territorial de equipamientos
universitarios. En Santiago, estos equipamientos están fuertemente concentrados en pocas comunas
centrales, por lo que esta variable también puede estar capturando centralidad urbana, concentración
de destinos educacionales y zonas con alta presencia de Metro o viajes asociados a \texttt{BIP}.

La densidad de juegos infantiles y comercio de conveniencia muestra asociaciones positivas con ambas
alternativas QR, especialmente con \texttt{QR\_RED}. Esto sugiere que ciertos entornos urbanos de
escala barrial se asocian con una mayor utilidad relativa de alternativas QR. La densidad de centros
deportivos presenta una señal menos clara: tiende a ser negativa para \texttt{QR\_OTHER} y no robusta
para \texttt{QR\_RED}.

La variable de refugios de transporte en OSM
(\path{osm_transport_shelter_yes_density_km2_z}) presenta una señal particularmente interesante:
el coeficiente es negativo y significativo para \texttt{QR\_OTHER}, pero positivo y significativo
para \texttt{QR\_RED}. Esto indica que la presencia de objetos de transporte con
\texttt{shelter=yes} no se asocia de forma homogénea con el uso de QR. Más que interpretarse como un
efecto causal de tener refugio, esta variable debe leerse como un indicador territorial de
infraestructura de parada o entorno de acceso al sistema.

La densidad de accesos a Metro
(\path{osm_railway_subway_entrance_density_km2_z}) presenta coeficientes negativos y significativos
para \texttt{QR\_OTHER} y \texttt{QR\_RED} respecto de \texttt{BIP}. Esto sugiere que zonas con
mayor presencia de accesos físicos a Metro se asocian relativamente menos con alternativas QR una vez
controladas las demás variables del modelo. La densidad de colegios también tiende a presentar una
señal negativa para \texttt{QR\_RED}, mientras que su asociación con \texttt{QR\_OTHER} no es
robusta.

\subsection{Variables EOD 2012}

Las variables de ingreso construidas desde la EOD 2012 se evalúan en dos formas alternativas. La
primera usa proporciones continuas de hogares en tramos históricos de ingreso; la segunda usa
indicadores de alta concentración zonal. Estas dos formas no responden exactamente la misma pregunta:
la especificación continua mide cambios graduales en la composición de ingreso de la zona, mientras
que la especificación dummy identifica zonas ubicadas en el cuartil superior de concentración de un
grupo específico.

En la especificación continua, la proporción histórica de hogares en los tramos bajos de ingreso
(\texttt{D+E\_proxy}) presenta una asociación positiva y significativa con \texttt{QR\_OTHER}, pero
no muestra una señal robusta para \texttt{QR\_RED}. En cambio, la proporción histórica de hogares en
el tramo superior de ingreso (\texttt{ABC1\_proxy}) presenta una señal negativa solo marginal para
\texttt{QR\_OTHER} y no robusta para \texttt{QR\_RED}. Esta lectura sugiere que el gradiente
continuo de ingreso bajo se relaciona principalmente con \texttt{QR\_OTHER}.

En la especificación dummy, la alta concentración de hogares \texttt{ABC1\_proxy} muestra
coeficientes negativos y significativos para \texttt{QR\_OTHER} y \texttt{QR\_RED}. En cambio, la
alta concentración de hogares \texttt{D+E\_proxy} no presenta una señal estadísticamente robusta.
Esto sugiere que, cuando el ingreso histórico se representa como una condición de alta concentración,
la señal más clara aparece en zonas de mayor ingreso relativo, que se asocian con menor utilidad
relativa de alternativas QR respecto de \texttt{BIP}.

Ambas lecturas son complementarias más que contradictorias. El modelo continuo permite observar una
relación gradual asociada a la presencia de hogares de ingresos bajos, especialmente para
\texttt{QR\_OTHER}; el modelo dummy, en cambio, resalta el comportamiento de zonas con concentración
alta de hogares de mayores ingresos. En cualquier caso, estas variables deben tratarse como proxies
históricos y territoriales, no como clasificación socioeconómica oficial ni como ingreso individual
del usuario.

\subsection{Síntesis}

En conjunto, los resultados muestran que la elección entre \texttt{BIP}, \texttt{QR\_OTHER} y
\texttt{QR\_RED} no solo depende de atributos operacionales del viaje, sino también de
características temporales y territoriales del origen. Las variables de espera y transbordo mantienen
los signos esperados en la mayoría de los casos, con una excepción sustantivamente interesante en
\texttt{QR\_RED}: la espera por transbordo no presenta una penalización estadísticamente robusta.

Los resultados más robustos para \texttt{QR\_RED} aparecen en la asociación positiva con educación
universitaria, puntas laborales, asistencia parvularia, proporción de mujeres y ciertos entornos OSM
de escala barrial. Para \texttt{QR\_OTHER}, destacan la asociación positiva con población inmigrante
y con el proxy continuo de hogares en tramos bajos de ingreso, junto con señales negativas asociadas
a hacinamiento, demanda local y refugios de transporte.

La comparación entre los modelos continuo y dummy muestra que la incorporación de la EOD 2012 no
altera los patrones principales del modelo, pero sí permite explorar dos formas distintas de
representar la dimensión socioeconómica histórica del territorio. Por lo tanto, la decisión entre
ambas especificaciones debería basarse principalmente en la interpretación sustantiva que se quiera
privilegiar: gradientes territoriales continuos de ingreso o identificación de zonas con alta
concentración relativa de ciertos grupos.
